\section{Optimization of an isolated rank-one flow block}\label{app:costs}

\subsection{Exact optimization of low-interaction-rank pool costs}
\label{s6:low-rank}
Consider the isolated rank-one margin block
\begin{equation}
 \mathcal K=\{W\in\R_+^{m\times n}:\rank W\le1,\quad
 \ell\le W\mathbf1\le u,\quad
 \ell'\le W^T\mathbf1\le u'\},\label{s6:margin-block}
\end{equation}
with finite nonnegative rational bounds in the stated order, and objective
$\langle C,W\rangle$. For $W\ne0$, its margins $x,y$ and common total
$S$ give $W=xy^T/S$ by \eqref{s1:rank-one-form}.
The feasible-total interval is
$I=[\max(\mathbf1^T\ell,\mathbf1^T\ell'),
\min(\mathbf1^Tu,\mathbf1^Tu')]$; reject an empty interval.
Define the \emph{interaction rank}
\[
 k=\rank D,\qquad D_{ij}=C_{ij}-C_{i1}-C_{1j}+C_{11}.
\]
This is the least rank after subtracting arbitrary additive row and
column costs. Indeed multiplying $C-a\mathbf1^T-\mathbf1b^T$ on the
left by $I_m-\mathbf1e_1^T$ and on the right by $I_n-e_1\mathbf1^T$
gives $D$, so no subtraction can have smaller rank. Taking
$a_i=C_{i1}$, $b_j=C_{1j}-C_{11}$ attains it.
Rational elimination therefore constructs
$C=a\mathbf1^T+\mathbf1b^T+UV^T$ with $k$ factor columns.
The objective is
\begin{equation}
 a^Tx+b^Ty+\frac{(U^Tx)^T(V^Ty)}S.\label{s6:interaction-objective}
\end{equation}

\begin{theorem}[Fixed interaction rank]\label{s6:rank-theorem}
For each fixed $k$, exact linear optimization over \eqref{s6:margin-block}
is polynomial in the rational input bit length, with dimension dependence
$(m+n)^{O(k+1)}$. An optimizer uses rational numbers and at most one
square root, all with polynomial encoding length. This is polynomial
for fixed $k$; the displayed dependence is not an FPT bound in $k$.
\end{theorem}
\begin{proof}
A linear image of an $h$-coordinate box in dimension $d$ is a translated
sum of $h$ segments. Its vertices correspond to full-dimensional sign
regions of the central hyperplanes orthogonal to its nonzero generators.
The region recurrence obtained by adding one hyperplane gives at most
$2\sum_{i=0}^{d-1}\binom{h-1}{i}$ vertices for positive effective
dimension, and one in dimension zero. Degenerate or repeated generators
only decrease this number. For fixed $d$, enumerate these vertices
constructively by adding segments one at a time: union the existing
vertices with their translates and retain only vertices of the resulting
convex hull. Membership of one candidate in the convex hull of the others
is a rational LP test. Store with every retained candidate a corner of
the original box mapping to it. Intermediate lists have the same
polynomial bound, and every stored coordinate is a rational generator
sum of polynomial bit length.

Apply this construction to
\[
 Z_x=\{(\mathbf1^Tx,U^Tx,a^Tx):\ell\le x\le u\},\qquad
 Z_y=\{(\mathbf1^Ty,V^Ty,b^Ty):\ell'\le y\le u'\}.
\]
Both projected boxes have dimension at most $k+2$ and polynomially many
vertices for fixed $k$. At any fixed $S>0$, the objective
\eqref{s6:interaction-objective} is affine in either projected margin
when the other is fixed. Starting with an optimum, minimize first on
one slice at total $S$ and choose a slice vertex, then do the same on
the other. Both steps preserve optimality. There is therefore an optimum
at a pair of slice vertices.

Every vertex of a hyperplane slice of a polytope lies on an original
vertex or edge. To see this, take its smallest containing face. The
point is in the relative interior of that face. If the face dimension
were at least two, intersecting with one hyperplane through the point
would leave a nonzero two-sided feasible direction. This contradicts
extremality in the slice.
Enumerate all pairs of projected-box vertices, including every edge.
A segment with different endpoint totals parametrizes a projected margin
affine in $S$ on the closed interval between those totals. Interpolating
its two stored box corners gives a feasible original-margin witness.
The additional nonedge segments are harmless, since they stay in the
projected box. Include each original vertex as a singleton candidate
at its fixed total. Equal-total segments can be omitted: an interior
point of an edge parallel to the slicing hyperplane cannot be a slice
vertex, and its endpoints are already included.

Pair row and column candidates and intersect their intervals with $I$.
On every positive interval, their projected coordinates and witnesses
are affine in $S$, and substitution gives
\begin{equation}
 f(S)=\alpha S+\beta+\gamma/S\label{s6:fractional-piece}
\end{equation}
with rational coefficients of polynomial bit length. Evaluate positive
endpoints and the stationary point $\sqrt{\gamma/\alpha}$ if
$\alpha,\gamma>0$ and it lies in the interval. These are all interior
minimum possibilities, apart from a constant function, where any
feasible total suffices. When both coefficients are negative a stationary
point is a maximum; the other sign cases are monotone. Singletons are
evaluated directly.

If $0\in I$, zero is feasible and is included separately; if $I=\{0\}$
it is the sole matrix. No missing infimum at zero occurs, since
$|\langle C,W\rangle|\le\|C\|_\infty S$ for every nonnegative
matrix of total $S$. In particular an interval candidate reaching zero
has the correct continuous zero extension. Every enumerated candidate
is feasible, and every fixed-total optimum has representatives on the
list, so minimizing the list gives a global optimum.
There are polynomially many candidates for fixed $k$. Interior values
are $\beta+2\sqrt{\alpha\gamma}$; comparisons of two such values
require sign checks and at most two squarings, or degree-four algebraic
comparison. Interpolation at the chosen total and $W=xy^T/S$ stay in
one quadratic number field. These operations and all rational LP tests
have polynomial bit complexity, proving the stated bound and output
representation.
\end{proof}

Low-rank box optimization and projected-box methods have established
precedents. Punnen, Sripratak and Karapetyan
\cite[Sections~3.2--3.3]{s6:punnen2015} treat bipartite box objectives;
Hlad\'ik, \v{C}ern\'y and Rada \cite{s6:hladik2021} use zonotope faces
for continuous quadratic box optimization. The additional feature here
is the shared variable total and its division in
\eqref{s6:interaction-objective}; the proof handles it by slice segments
and one-variable rational optimization.

\begin{proposition}[Rank-one cost specialization]\label{s6:rank-one-cost}
If $C=pq^T$ is supplied in factored form, the optimum and its margins
can be found in $O((m+n)\log(m+n))$ arithmetic operations and comparisons,
with polynomial bit complexity. Factoring a supplied dense rank-one
matrix or writing every entry of $W$ adds $O(mn)$ arithmetic work.
\end{proposition}
\begin{proof}
At each positive total define $A_-(S),A_+(S)$ as the minimum and maximum
of $p^Tx$ on $[\ell,u]\cap\{\mathbf1^Tx=S\}$, and define $B_-,B_+$
similarly. The two scalar images are intervals, so their product is
minimized at one of the four endpoint pairs, regardless of signs.
The optimal value is
$\min_{\sigma,\tau\in\{-,+\}}A_\sigma(S)B_\tau(S)/S$.
Each endpoint function is continuous knapsack: start at the lower bounds
and fill residual capacities in increasing coefficient order for a
minimum, decreasing order for a maximum. Sorting and prefix sums give
all affine pieces and $O(m+n)$ total breakpoints. Merge those breakpoints
with $I$ and sweep the four paired functions. Each is of form
\eqref{s6:fractional-piece}; process its endpoints and possible minimum
stationary point as above. Recover the selected greedy margins in linear
work. The same exact comparisons prove polynomial bit complexity.
\end{proof}
For example, fix the first row and column margins to one, bound the
other two margins by $[0,1]$, and take $p=(3,1)$, $q=(2,1)$.
The cost is $S+3+2/S$ for $1\le S\le2$, minimized at $S=\sqrt2$.
Even a rank-one cost need not have a rational optimizer.

\begin{corollary}[A fixed-attribute Lagrangian pool oracle]\label{s6:oracle}
Suppose an isolated pool block has additive operating costs, row and
column margin bounds as in \eqref{s6:margin-block}, and $K$ attributes.
For any supplied rational nonnegative multipliers on terminal quality
inequalities, the Lagrangian subproblem obtained by dualizing those
inequalities is exactly solvable in polynomial time for fixed $K$.
\end{corollary}
\begin{proof}
Writing input attributes $Q_{ik}$ and output upper bounds $H_{jk}$,
the homogeneous quality rows are
$\sum_i(Q_{ik}-H_{jk})W_{ij}\le0$.
For minimization, multipliers $\lambda_{jk}\ge0$ add matrix cost
\[
 Q\Lambda^T-\mathbf1 h^T,\qquad h_j=\sum_k\lambda_{jk}H_{jk}.
\]
The additive part does not affect interaction rank, and
$\rank(Q\Lambda^T)\le K$. Apply \cref{s6:rank-theorem}.
Lower quality rows reverse the corresponding signs; both upper and lower
multipliers can be combined into the same $K$ input-attribute columns.
An affine basis of the input qualities can further replace $K$ by its
rank. All remaining constraints must have precisely the margin/rank form;
other pool couplings need separate treatment. This proves an exact
subproblem oracle, with no zero-duality-gap or full pooling tractability
claim.
\end{proof}
If a supplied cost approximation $\widetilde C$ has fixed interaction
rank and $\|C-\widetilde C\|_\infty\le\eta$, where this is the maximum
absolute entry norm, its optimizer has true cost at most
$2\eta S_{\max}$ above optimum, with
$S_{\max}=\min(\mathbf1^Tu,\mathbf1^Tu')$. Indeed every feasible
matrix has objective perturbation at most $\eta S_{\max}$. The bound
assumes both the approximation and its error certificate are supplied.


\subsection{General costs: hardness and short exact witnesses}

The fixed-interaction-rank algorithm does not extend to arbitrary costs
unless $\PP=\NP$. We prove this for the same margin set
\eqref{s6:margin-block}, including the zero-lower-bound restriction.
The reduction answers the precise conjecture after Theorem~4 of
\citet{dey2020-convexifications-of-rank-one-based}. It uses the
established NP-completeness of maximum edge biclique
\cite{app:peeters2003}; no priority is claimed for its underlying
Boolean bilinear encoding. \citet{app:gillis2014} already connect
maximum edge biclique with nonnegative rank-one approximation of a
possibly signed matrix under a squared Frobenius objective. Here the
objective is linear and the constraints are simultaneous margin bounds.

\begin{theorem}[General-cost margin hardness]\label{app:margin-hardness}
Deciding whether $\min_{W\in\mathcal K}\langle C,W\rangle\le0$ is
strongly $\NP$-hard, even with all margin bounds in $\{0,1\}$ and
integer costs of polynomially bounded magnitude. The only positive
lower bounds needed are one row margin and one column margin, both one.
\end{theorem}
\begin{proof}
Given a bipartite graph with sides of sizes $p,q$, let
$h=\max(p,q)$ and let $1\le k\le|E|$ be the required biclique edge
count. Use one distinguished row and column, each of exact margin one;
$p$ graph rows and $q$ graph columns; and $h$ dummy rows and $h$ dummy
columns. Every remaining margin has interval $[0,1]$. Set $C_{00}=k$,
set graph-block entries to $-1$ on edges and $h^2$ on nonedges, and
set every other entry to zero. Writing $x,y$ for graph margins,
the numerator in \eqref{s1:rank-one-form} is
\[
 r^TCc=k+\sum_{i,j}C_{ij}x_i y_j=k+Q(x,y).
\]
Dummy margins enforce equal totals for every
$(x,y)\in[0,1]^p\times[0,1]^q$: their required total difference lies
in $[-h,h]$, and each dummy side realizes every total in $[0,h]$.
The common total $S$ is at least one, so the objective is nonpositive
exactly when this numerator is nonpositive.

Minimize the bilinear function $Q$ first over the row box with $y$
fixed, choosing a binary minimizer, and then over the column box with
that row choice fixed. Its minimum therefore occurs at binary $x,y$.
For their supports $A,B$, a selected nonedge gives
$Q\ge h^2-(|A||B|-1)\ge1$. If there is no nonedge, then
$k+Q=k-|A||B|$. Thus the constructed threshold holds exactly when
there is a biclique with at least $k$ edges. Dimensions are linear in
$h$ and all numbers are bounded by $h^2$, proving strong hardness.
In a no instance the numerator is at least one everywhere, because
its minimum occurs at a binary pair. Since $S\le2h+1$, the optimal
value is then at least $1/(2h+1)$.
\end{proof}

\begin{lemma}[Repair of two unit margins]\label{app:unit-repair}
Let $\mathcal U_{p,q}$ be the nonnegative rank-one matrices whose row
and column margins are at most one, and let
$F=\{W\in\mathcal U_{p,q}:r_0=c_0=1\}$. Every
$W\in\mathcal U_{p,q}$ has $W'\in F$ satisfying
\[
 \|W-W'\|_1\le2(2-r_0-c_0),
 \qquad \|W\|_1=\sum_{i,j}|W_{ij}|.
\]
\end{lemma}
\begin{proof}
Put $\delta=2-r_0-c_0$. When $S\ge1$, raise $r_0$ to one and
reduce the other row margins by $1-r_0$, preserving total $S$;
their total $S-r_0$ suffices. Do the same to the column margins.
The repaired vectors $r',c'$ remain in their unit boxes and give
$W'=r'(c')^T/S$. The outer-product triangle inequality gives
\[
 \|W'-W\|_1\le\|r'-r\|_1+\|c'-c\|_1=2\delta.
\]
When $0<S<1$, take $W'=E_{00}$. From
$(S-r_0)(S-c_0)\ge0$ follows $W_{00}\ge r_0+c_0-S$, and hence
\[
 \|W-E_{00}\|_1=1+S-2W_{00}
 \le1+3S-2r_0-2c_0\le2\delta.
\]
At $S=0$ the same choice satisfies the bound directly.
\end{proof}

\begin{corollary}[Unit upper bounds and zero lower bounds]
\label{app:unit-hardness}
Rational threshold optimization over $\mathcal U_{p,q}$ is strongly
$\NP$-hard with integer costs and a negative threshold, all of
polynomially bounded magnitude.
\end{corollary}
\begin{proof}
For the matrix $C$ in \cref{app:margin-hardness}, put
$B=\max_{ij}|C_{ij}|$ and $L=2B+1$. Remove the two positive lower
bounds and replace the objective by
$\langle C,W\rangle-L(r_0+c_0)$, a linear matrix cost.
For a point outside $F$, the repair in \cref{app:unit-repair} changes
this objective by at most $(2B-L)\delta<0$. Every optimizer therefore
belongs to $F$, and the optimum is exactly its previous value minus
$2L$. Use threshold $-2L$. New cost magnitudes are at most $5B+2$;
the construction is strong and retains the gap $1/(2h+1)$.
\end{proof}

\begin{theorem}[Margin certificates and algebraic optima]
\label{app:margin-certificates}
For finite rational data in \eqref{s6:margin-block}, every nonempty
instance has an optimizer in one real quadratic number field with
polynomial encoding length. Every yes instance for a rational cost
threshold has a rational matrix witness of polynomial bit length.
Consequently rational threshold decision is strongly $\NP$-complete.
\end{theorem}
\begin{proof}
Compactness gives an optimum. Treat the feasible zero matrix separately.
At an optimal positive total $S_*$, minimize first over the row-margin
slice with the column margins fixed, then over the column slice with
those rows fixed. Choose a vertex at each step; both steps preserve
global optimality. Every vertex of a box intersected with one total
hyperplane has all but at most one coordinate at a bound: two strictly
interior coordinates would allow opposite perturbations. Fix these two
margin patterns. They are affine functions
\[
 r(S)=e_aS+b,\qquad c(S)=e_hS+d
\]
on a closed rational interval $J$, with all coefficients of polynomial
bit length. Here $e_a,e_h$ denote unit vectors in their respective spaces.
Their objective for $S>0$ is
$f(S)=\alpha S+\beta+\gamma/S$, where
\[
 \alpha=C_{ah},\quad
 \beta=e_a^TCd+b^TCe_h,\quad \gamma=b^TCd.
\]
These expressions have polynomial bit length. The endpoint and
stationary-point argument of \cref{s6:rank-theorem} shows that a global
optimum in this family is rational or obtained at
$S=\sqrt{\gamma/\alpha}$ with $\alpha,\gamma>0$. Its margins and
matrix belong to the same quadratic field. If $J$ reaches zero,
nonnegative margins with total zero vanish, and the continuous
zero extension is handled separately.

Now let the rational threshold be $t$. The family contains a positive
threshold witness exactly when, for some positive $S\in J$,
\[
 q(S)=\alpha S^2+(\beta-t)S+\gamma\le0.
\]
A quadratic achieves its minimum on a rational closed interval at a
rational endpoint or, for positive leading coefficient, at the rational
stationary point $(t-\beta)/(2\alpha)$. If the minimum is negative,
its minimizer cannot be zero: if $0\in J$, then $r(0)=c(0)=0$, so
$q(0)=0$. If the minimum equals zero, any positive threshold witness is
a minimizer. Such a minimizer is a positive rational endpoint or the
rational stationary point, unless $q$ is constant; in the constant case
take any positive rational point of $J$. This gives a rational total,
margins, and matrix of polynomial bit length.
A verifier checks margin bounds, rank at most one, and the rational
objective inequality in polynomial time. Together with
\cref{app:margin-hardness,app:unit-hardness}, this proves the last claim.
\end{proof}

\begin{theorem}[One fixed margin dimension]\label{app:margin-fpt}
For $d=\min(m,n)$ and total rational input length $B$, exact optimization
over \eqref{s6:margin-block} takes
$d2^d\poly(m,n,B)$ bit operations. In particular it is polynomial
when either matrix dimension is fixed.
\end{theorem}
\begin{proof}
Transpose so the row dimension is $d$. For any fixed positive total,
minimizing over column margins makes the objective numerator a concave
function of row margins. A minimum occurs at a row-slice vertex.
Enumerate its free coordinate and a lower/upper choice for every other
coordinate, at most $d2^{d-1}$ patterns. Each pattern gives an affine
row vector $r(S)$ and a rational validity interval. The column weights
$(C^Tr(S))_j$ are affine in $S$. Their $O(n^2)$ pairwise crossings split
the interval into pieces with a fixed weight order; identical weights
use a fixed index order. On each piece, the continuous-knapsack rule
fills column margins from their lower bounds in that order. Its prefix
capacities contribute at most $n$ further breakpoints. Both margins
are therefore affine on $O(n^3)$ intervals per row pattern.

At ties either adjacent greedy order is optimal, and singleton intervals
are checked as well. On every positive interval apply the endpoint and
square-root calculation for \eqref{s6:fractional-piece}; include zero
when feasible. Every candidate is feasible and the enumeration covers
an optimum at every total. All rational arithmetic and algebraic
comparisons have polynomial bit complexity as proved in
\cref{s6:rank-theorem}. The factor outside the polynomial is precisely
the row-pattern count, giving the claimed dependence.
\end{proof}

For example, fix the first row and column margins to one, let the other
margins vary in $[0,1]$, and use $C=I_2$. Equal totals force the remaining
margins to be $S-1$, and the objective is $S-2+2/S$ on $[1,2]$.
Its unique minimum occurs at $S=\sqrt2$. Thus rational threshold
certificates do not imply rational exact optimizers.
